% C_6 and 2C_3 share the trivial 1-block fair partition and quotient matrix (2),
% but have different residual spectra on W^\perp.
\begin{tikzpicture}[x=1mm,y=1mm,font=\figurefont\fontsize{8}{10}\selectfont,
  vertex/.style={circle,draw=FigureBlue,fill=FigureBlueFill,inner sep=0pt,
    minimum size=3.4mm,line width=.55pt},
  edge/.style={draw=FigureStroke,line width=.65pt},
  clsbox/.style={draw=FigureBlue,dashed,fill=FigureBlueFill!35!white,rounded corners=2pt,line width=.55pt},
  note/.style={font=\figurefont\fontsize{7.5}{9}\selectfont,text=FigureInk},
  panel/.style={font=\figurefont\fontsize{7.5}{9}\selectfont,text=FigureInk}]
  % Panel (a): C_6
  \begin{scope}[xshift=0mm]
    \filldraw[clsbox] (-14,-12.5) rectangle (14,12.5);
    \foreach \i in {0,...,5} {
      \node[vertex] (h\i) at ({90+60*\i}:9) {};
    }
    \draw[edge] (h0)--(h1)--(h2)--(h3)--(h4)--(h5)--(h0);
    \node[note,anchor=north west] at (-13,12) {$C_1$};
    \node[note] at (0,-16.5) {$\bm Q=(2),\ \text{商谱 }\{2\}$};
    \node[note,text=FigureMuted] at (0,-21) {$\text{剩余谱：}1,1,-1,-1,-2$};
    \node[panel] at (0,-27) {(a) 六点环 $C_6$};
  \end{scope}

  % Panel (b): 2C_3
  \begin{scope}[xshift=48mm]
    \filldraw[clsbox] (-18,-12.5) rectangle (18,12.5);
    \node[vertex] (t1) at (-9,6.5) {};
    \node[vertex] (t2) at (-14,-4.5) {};
    \node[vertex] (t3) at (-4,-4.5) {};
    \draw[edge] (t1)--(t2)--(t3)--(t1);
    \node[vertex] (s1) at (9,6.5) {};
    \node[vertex] (s2) at (4,-4.5) {};
    \node[vertex] (s3) at (14,-4.5) {};
    \draw[edge] (s1)--(s2)--(s3)--(s1);
    \node[note,anchor=north west] at (-17,12) {$C_1$};
    \node[note] at (0,-16.5) {$\bm Q=(2),\ \text{商谱 }\{2\}$};
    \node[note,text=FigureMuted] at (0,-21) {$\text{剩余谱：}2,-1,-1,-1,-1$};
    \node[panel] at (0,-27) {(b) 两个三角形 $2C_3$};
  \end{scope}
\end{tikzpicture}
